A partir de los resultados obtenidos, se concluye que en los
algoritmos de multiplicación de matrices, Naive resulta más
eficiente que Strassen en tiempo de ejecución para todas las
dimensiones evaluadas hasta $n = 1024$, evidenciando de que la
sobrecarga de submatrices intermedias de Strassen ralentiza su
desempeño y genera además un consumo de memoria notablemente mayor.
En el ordenamiento de arreglos, el algoritmo Sort se consolida como
la alternativa más robusta y eficiente, tanto en tiempo de ejecución
como en uso de memoria auxiliar. También se demostro la alta
sensibilidad que tiene el algoritmo QuickSort con partición de Lomuto
ante entradas ordenados, donde su tiempo de ejecución empeora drásticamente
de $O(n \log n)$ a $O(n^2)$, mientras que MergeSort garantiza
estabilidad temporal ante cualquier escenario de datos a costa de
un costo espacial estricto $O(n)$ frente a las soluciones in-place.